% Title (docs/PAPER_OUTLINE.md): Label granularity in supervised pretraining of jet foundation models
% Abstract text only; main.tex supplies the environment. PI brief for draft 94: about 200 words, Sophon's
% shape (context, gap, what we did, three findings, recommendation), no definitions, no numbers, Sophon
% expanded at first use.
% CHECK: the findings sentences carry the same CHECKs as the last paragraph of the Introduction
Jet foundation models are often pretrained to classify simulated jets into many fine classes, and analyses
build discriminants by summing the scores of those classes. Whether the granularity of these labels
changes what the network learns has not been tested with the training jets held fixed. We pretrain the
network of Signature-Oriented Pre-training for Heavy-resonance ObservatioN (Sophon) repeatedly on the same
jets, in the same order and from the same initial weights, and change only the label set, along a
hierarchy from Sophon's full set of classes to one that describes only the number and type of the prongs.
Random partitions of the classes, some differing only in a flavor boundary, compare label sets at a fixed
number of classes. Coarser labels make linear probes on the frozen features substantially worse on the
pretraining classes. Merging quark flavor is the costliest single merge, and at a fixed number of classes, which classes
are merged changes the cost. After fine-tuning, the coarsest labels cost a few per cent more errors on
JetClass and no resolved difference on the public top-tagging and quark/gluon benchmarks. An auxiliary
jet-mass regression makes the mass predictable from the frozen features. A pretraining
label set should contain every distinction that downstream discriminants sum over, and the flavor
boundaries.
